\documentclass{article}
\usepackage[utf8]{inputenc}
\usepackage[T1]{fontenc}
\usepackage[margin=2cm]{geometry}
\usepackage{hyperref}
\sloppy
\begin{document}
\section*{Historical source}
Encyclopaedia Britannica scan, source row 1284. Parse 3.5 extracted the text. Super 3.5 VL corrected it from an enlarged crop of each block. It can still contain recognition errors.
Page headers, footers, and page numbers are removed.

\subsection*{Untitled continuation (E0106)}
\textit{Continues an entry from an earlier page.}

Scotland. the meanest of the human race. To accuse me of felonying the death of the queen my sister, is to brand me with the infamy which I abhor most. It is my nature to employ the devotions of Esther, and not the sword of Judith. Elizabeth herself will attest, that I have often admonished her not to draw upon her head the resentment of my friends by the enormity of her cruelties to me. My innocence cannot sincerely be doubted; and it is known to the Almighty God, that I could not possibly think to forego his mercy, and to ruin my soul, in order to compass a transgression so horrible as that of her murder. But amidst the inclement and unprincipled pretences which my adversaries are pleased to invent to overwhelm me with calamities and anguish, I can trace and discover with ease the real causes of their hostility and provocation. My crimes are my birth, the injuries I have been compelled to endure, and my religion. I am proud of the first; I can forgive the second; and the third is a source to me of such comfort and hope, that for its glory I will be contented that my blood shall flow upon the scaffold.

To the defence of Mary, no returns were made before the trial and unsupported affirmations of the truth of the evidence produced to her prejudice. In the course of the trial, however, there occurred some incidents which deserve to be related. My lord Burleigh, who was willing to disfcompone her, charged her with a fixed resolution of conveying her claims and titles to England to the king of Spain. But though, in a discontented humour with her son, he had threatened to disinherit him, and had even corresponded on the subject with her select friends, it appears that this project is to be considered as only a transient effect of resentment and passion. She indeed acknowledged, that the Spaniard professed to have pretensions to the kingdom of England, and that a book in justification of them had been communicated to her. She declared, however, that she had incurred the displeasure of many by disapproving of this book; and that no conveyance of her titles to the Spaniard had been ever executed.

The trial continued during the space of two days; but the commissioners avoided to deliver their opinions. My lord Burleigh, in whose management Elizabeth chiefly confided, and whom the Scottish queen discomposed in no common degree by her ability and vigour, being eager to conclude the business, demanded to know if she had any thing to add to what she had urged in her defence. She informed him, that she would be infinitely pleased and gratified, if it should be permitted to her to be heard in her justification before a full meeting of the parliament, or before the queen and her privy-council. This intimation was unexpected; and the request implied in it was rejected. The court, in consequence of previous instructions from Elizabeth, adjourned to a further day, and appointed that the place of its convention should be the star-chamber at Westminster. It accordingly assembled there; and Naw and Curl, who had not been produced at Fotheringay castle, were now called before the commissioners. An oath to declare the truth was put to them; and they definitely affirmed and protested that the declaration they had subscribed was in every respect just and faithful. Nothing farther remained,

but to pronounce sentence against Mary. The commissioners unanimously concurred in delivering it as their verdict or judgment, that she "was a party to the conspiracy of Babington; and that she had compassed and imagined matters within the realm of England tending to the hurt, death, and destruction of the royal person of Elizabeth, in opposition to the statute framed for her protection." Upon the same day in which this extraordinary sentence was given, the commissioners and the judges of England issued a declaration which imported, that it was not to derogate in any degree from the titles and honour of the king of Scots.

The sentence against Mary was very soon afterwards ratified by the English parliament. King James was struck with horror at hearing of the execution of his mother; but that spirited prince could show his resentment no farther than by unavailing embassies and remonstrances. France interposed in the same ineffectual manner; and on the 6th of December 1586, Elizabeth caused the sentence of the commissioners against her to be proclaimed. After this she was made acquainted with her fate, and received the news with the greatest composure, and even apparent satisfaction. Her keepers now refused to treat her with any reverence or respect. They entered her apartment with their heads covered, and made no obeisance to her. They took down her canopy of state, and deprived her of all the badges of royalty. By these insulting mortifications they meant to inform her, that she had sunk from the dignity of a princess to the abject state of a criminal. She smiled, and said, "In despite of your sovereign and her subservient judges, I will live and die a queen. My royal character is indelible; and I will surrender it with my spirit to the Almighty God, from whom I received it, and to whom my honour and my innocence are fully known." In this melancholy situation Mary addressed a magnanimous letter to Elizabeth, in which, without making the least solicitation for her life, she only requested that her body might be carried to France; that she might be publicly executed; that her servants might be permitted to depart out of England unmolested, and enjoy the legacies which she bequeathed them. But to this letter no answer was given.

In the mean time James, who had neither address nor courage to attempt any thing in behalf of his mother, announced her situation to his bigotted subjects, and ordered prayers to be said for her in all the churches. The form of the petition to be preferred was framed with delicacy and caution, that the clergy might have no objection to it. He enjoined them to pray, "that it might please God to enlighten Mary with the light of his truth, and to protect her from the danger which was hanging over her." His own chaplains, and Mr David Lindley minister of Leith, observed his command. But all the other clergy refused to prostitute their pulpits by preferring any petitions to the Almighty for a Papist. James, shocked with their spirit of intolerance and sedition, appointed a new day for prayers to be said for Mary, and issued a stricter injunction to the clergy to obey him; and that he might be free himself from any insult, he commanded the archbishop of St Andrews to preach before him. The ecclesiastics, disgraced with his injunction,

She desires to be heard before the parliament, or before the queen.


\section*{Source record}
\begin{description}
\item[Scan ID] \texttt{britannica:309f765b1ecc83e0a7f7762e44b4a927a7555443:train:1284}
\item[Image file] \texttt{page-1284-b6a1ee180d42.jpg}
\item[Image SHA-256] \texttt{b6a1ee180d4254f561bccaf072231163bd53aa804dad1fd5718590f53b90f394}
\item[Dataset] NationalLibraryOfScotland/encyclopaedia\_britannica\_illustrated at revision \texttt{309f765b1ecc83e0a7f7762e44b4a927a7555443} (CC0-1.0)
\item[Extraction] model \texttt{nvidia/nemotron-3.5-parse}, raw reply \texttt{extract/raw/page-1284-b6a1ee180d42.json}
\item[Raw reply SHA-256] \texttt{8eb64b5a30714843c0653da9315a05bd58d7bab9d177ccead9d77a5ca4b4d80e}
\item[Correction] model \texttt{nvidia/Nemotron-3.5-Super-VL}, block crops, raw replies \texttt{correct/raw/crop\_nothink/page-1284-b6a1ee180d42-b*.json}; super\_vl: 7
\item[Running header] S C O [ 865 ]
\end{description}
\end{document}
